\doublespacing % Do not change - required

\chapter{Introduction}
\label{ch1}

%%%%%%%%%%%%%%%%%%%%%%%%%%%%%%%%%%%%%%%
% IMPORTANT
\begin{spacing}{1}
\minitoc
\end{spacing}
\thesisspacing
%%%%%%%%%%%%%%%%%%%%%%%%%%%%%%%%%%%%%%%

\section{Motivation}
Machine learning systems are increasingly deployed across networks of devices that observe, decide, and interact with the physical world. In many deployments, training data is inherently distributed across mobile devices, embedded sensors, and edge gateways, and cannot be conveniently centralised due to privacy constraints, regulatory requirements, bandwidth limitations, or operational risk. Federated learning addresses this setting by enabling clients to compute model updates on local data and communicate only updates to a coordinating server.

However, decentralisation enlarges the threat surface. Standard aggregation rules, such as averaging, assume that received updates are broadly informative and aligned with the global learning objective. In realistic systems this assumption can be violated by device faults, software bugs, unreliable networks, or adversarial clients. Under such conditions, a small number of abnormal updates may dominate the aggregate, destabilise training, or steer the model towards unintended behaviour. Robust aggregation is therefore essential for trustworthy distributed learning.

Robustness is further complicated by client heterogeneity. Under non-identically distributed (non-IID) data, benign clients may legitimately differ in update direction and magnitude, and classical Byzantine-robust rules that rely on tight clustering can over-suppress these benign but diverse contributions. Finally, temporal dynamics introduce additional vulnerabilities: server-side momentum and other history-based mechanisms improve optimisation, but can be exploited by stealthy adversaries that apply small perturbations over many rounds, gradually biasing server state and inducing long-term drift. These challenges motivate history-aware and non-IID-tolerant aggregation methods that suppress malicious behaviour without relying on rigid thresholds or oracle knowledge of the attacker fraction.

\section{Project Aims and Contributions}
This project investigates robustness in federated learning under the Byzantine
threat model, where a subset of clients may transmit arbitrary or adversarial
updates. The aim is to design and evaluate a server-side aggregation mechanism
that remains useful under both IID and non-IID data while resisting attacks that
act through update magnitude, update structure, or persistent temporal drift.
The central difficulty is that benign heterogeneity and malicious behaviour can
look similar at the level of a single communication round: a non-IID client may
legitimately differ from the population, while a Byzantine client may exploit
that diversity to avoid simple outlier checks.

The main technical contribution is \emph{HRAC} (History-Residual Adaptive
Clipping), a history-aware aggregation rule that compares each client primarily
with its own past behaviour rather than with a single global centre. HRAC
combines a median-based global norm cap, per-client residual clipping, and
soft weighting based on residual-change statistics. This design is intended to
reduce the influence of harmful updates without turning the aggregator into a
hard detector that requires attacker labels, an oracle attacker fraction, or a
tight cluster of benign clients.

The project also contributes an implementation and evaluation pipeline for
studying this design in controlled federated learning experiments. The pipeline
supports interchangeable attacks and aggregation baselines, logs diagnostic
quantities such as residual thresholds, $\mu/\nu$ histories and client weights,
and produces accuracy curves and ablation results for analysing robustness--
utility trade-offs. These implementation and diagnostic components are part of
the contribution because they make it possible to assess not only whether HRAC
improves final accuracy in a given setting, but also which parts of the method
are responsible for its behaviour.

The remainder of the report is organised as follows. Chapter~\ref{ch2} reviews
federated learning, Byzantine-robust aggregation, poisoning attacks, and
history-aware defences. Chapter~\ref{ch:method-implementation} presents the
HRAC method and its implementation. Chapter~\ref{ch:experiments-results}
evaluates HRAC against baseline aggregators and reports component ablations.
The final chapters discuss professional issues, project management,
limitations, and future work.
